\section{LoRA}

Low-Rank Adaptation (LoRA) is a parameter-efficient\footnote{Parameter-efficient methods train relatively few task-specific parameters rather than updating every pre-trained model parameter.} fine-tuning\footnote{Fine-tuning adapts a pre-trained model to a downstream task using task-specific data.} method. It freezes pre-trained weights and learns low-rank updates for selected weight matrices. These notes follow \href{https://arxiv.org/abs/2106.09685}{Hu et al., \emph{LoRA: Low-Rank Adaptation of Large Language Models}}.

\subsection{Introduction}

\subsubsection{The problem}

Full fine-tuning updates all model parameters. For a large model such as GPT-3 (175 billion parameters), this makes training expensive and requires storing a separate full-size checkpoint for each task. Other parameter-efficient methods reduce task-specific storage, but some add inference latency or use part of the available context window.

\subsubsection{The solution}

Motivated by evidence that model adaptation can take place in a low-dimensional subspace, the LoRA authors hypothesize that the required weight updates have low intrinsic rank.\footnote{Here, low intrinsic rank means that the task-specific change to a weight matrix can be represented effectively by a product of two matrices with a small shared dimension.} LoRA therefore represents the update to selected weight matrices with two trainable, low-rank matrices while keeping the original weights fixed.

This gives LoRA several practical advantages:
\begin{itemize}
    \item A single pre-trained model can be shared across tasks; each task needs only its own small LoRA matrices.
    \item Training uses fewer trainable parameters and less memory for gradients and optimizer states than full fine-tuning.
    \item The update can be merged into the original weights for inference, so it adds no inference latency.\footnote{Inference latency is the time a model takes to produce an output after receiving an input.}
    \item LoRA can be combined with other adaptation methods.
\end{itemize}

\subsection{Formulation}

\subsubsection{Problem statement}

Consider a pre-trained autoregressive language model\footnote{An autoregressive language model predicts each next token conditioned on the input and the tokens generated so far.} $P_{\Phi}(y\mid x)$ with parameters $\Phi_0$. For a downstream task, let $\mathcal{Z}=\{(x_i,y_i)\}_{i=1}^{N}$ be context-target pairs, where $x_i$ and $y_i$ are token sequences. Full fine-tuning updates $\Phi_0$ to maximize the conditional log-likelihood:
\[
\max_{\Phi}
\sum_{(x,y)\in\mathcal{Z}}
\sum_{t=1}^{|y|}
\log P_{\Phi}(y_t\mid x,y_{<t}).
\]

This approach learns a full-size parameter change $\Delta\Phi$ for every task. LoRA instead learns a much smaller set of parameters $\Theta$, which defines updates $\Delta\Phi(\Theta)$ to selected weights while the pre-trained parameters stay frozen. The objective becomes
\[
\max_{\Theta}
\sum_{(x,y)\in\mathcal{Z}}
\sum_{t=1}^{|y|}
\log P_{\Phi_0+\Delta\Phi(\Theta)}(y_t\mid x,y_{<t}),
\qquad |\Theta|\ll|\Phi_0|.
\]

\subsubsection{Method}

For a pre-trained weight matrix $W_0\in\mathbb{R}^{d\times k}$, LoRA constrains its update to $\Delta W=BA$, where $B\in\mathbb{R}^{d\times r}$, $A\in\mathbb{R}^{r\times k}$, and $r\ll\min(d,k)$. With input $x\in\mathbb{R}^{k}$, the adapted layer computes
\[
h=W_0x+\frac{\alpha}{r}BAx,
\]
where $\alpha$ is a scaling hyperparameter. The factor $\alpha/r$ reduces the need to retune other hyperparameters when the rank $r$ changes. Only $A$ and $B$ are trained; $W_0$ remains frozen.

For this matrix, LoRA trains $r(d+k)$ parameters instead of the $dk$ weights that full fine-tuning would update. This reduces the task-specific checkpoint size and the gradient and optimizer-state memory for the adapted matrix. The base weights must still be stored and used.

\subsubsection{Initialization and deployment}

The original LoRA method initializes $A$ with random Gaussian values and $B$ with zeros. Consequently, $BA=0$ at the start of training, so the adapted layer initially computes exactly $W_0x$. Training then changes $A$ and $B$ while leaving $W_0$ fixed.

After training, the update can be merged into the weight matrix:
\[
W_{\mathrm{merged}}=W_0+\frac{\alpha}{r}BA.
\]
Inference with $W_{\mathrm{merged}}x$ uses the same matrix multiplication as the original layer and incurs no additional latency from a separate LoRA path. For multiple tasks, one can retain the shared base model and store a small adapter for each task, merging the chosen adapter when needed. Deploying an adapter separately still requires access to its matching base model.
